\section{Cloud Microphysics}


Cloud processes include:


\begin{itemize}

\item Cloud water

\item Cloud ice

\item Rain

\item Snow

\item Graupel

\item Hail

\end{itemize}



The hydrometeor mixing ratios are represented by:


\begin{equation}
q_x
\end{equation}


where $x$ represents each hydrometeor category
($x \in \{v,c,r,i,s,g\}$: vapor, cloud water, rain, cloud ice, snow,
graupel). Each category obeys a flux-form conservation law with a
microphysical source/sink term $S_x$:

\begin{equation}
\frac{\partial(\rho q_x)}{\partial t}
+
\nabla\cdot(\rho q_x \vec{V})
=
S_x
\end{equation}


The dominant warm-rain source terms follow the Kessler (1969)
parameterization. Autoconversion of cloud water to rain past a critical
threshold $q_{c,crit}$:

\begin{equation}
P_{auto} = k_{auto}\max(q_c - q_{c,crit},\,0)
\end{equation}


Collision-coalescence growth of rain by collection of cloud water:

\begin{equation}
P_{coll} = k_{coll}\, q_c\, q_r^{7/8}
\end{equation}
